\scanpage{44}
% archive sheet number: 43

\textcircled{\scriptsize 15} \uline{Corps alg. clos} $=$ corps ultralocal \quad
$X_{\rm corps\,alg\,clos}\simeq\widetilde R_{\rm top\,corps\,alg\,clos}$

topologie où $X_i\to X$ est couvrante si elle est surjective
$\simeq\widetilde R_{\mathrm{cons\,fppf}}$\quad $R_{\rm cons}$ comme dans 11)
topologie fppf $=$ topologie des familles \deleted{\unsure{contenant uniquement des|contenant une}} familles \deleted{surjectives} finies surjectives.

\textcircled{\scriptsize 16} \uline{Anneau \deleted{\unsure{en}} annulé par $n$} ($\geq1$)

\[
 \text{\textcircled{\scriptsize a}}\qquad
 \begin{aligned}n1_A&=0\\ \text{ou}\qquad V(n1_A)&\lhook\joinrel\xrightarrow[\rm isom]{} e_{\mathcal E}\end{aligned}
\]

\[
 X_{\deleted{n\text{-tors}}\mathbb Z/n\mathbb Z}\simeq\widetilde R_{\mathrm{top}\ n\text{-tors}}
 \simeq \widehat R_{\mathbb Z/n\mathbb Z}
\]

% slanted
\margin{\deleted{\unsure{X est distinct}}\quad Sous-topos ouvert de $X_{\rm ann}=\widehat R$\quad
\uline{N.B} $X_{\unsure{1\text{-tors}|1\text{-tors.}}}=X_\varnothing=$ topos ponctuel}

top $n$-tors topologie où $X_i\to X$ est couvrante ssi un des $V(n1_{X_i})\to V(n1_X)$ \deleted{\unsure{est couvrant}} \inserted{une} section
($X_i\times_{\mathbb Z}\mathbb Z/\deleted{\unsure{1}}n\to X\times_{\mathbb Z}\mathbb Z/n\mathbb Z$)
$R_{\mathbb Z/n\mathbb Z}$ \deleted{\unsure{top. de}}\inserted{sous cat.} ple de $R$ formée des \deleted{\unsure{ens}} $X$ qui sont des $\mathbb Z/n\mathbb Z$-schémas

($R_{\mathbb Z/n\mathbb Z}$ peut être considéré comme une catégorie de fractions de $R$ pour les flèches $f$ telles que $f\otimes_{\mathbb Z}\mathbb Z/n\mathbb Z$ bijections)

\deleted{\unsure{b)}} \uline{Ex} Si $n$ premier, ce sont les ``anneaux de car $p>0$

\deleted{\unsure{a)}} \textcircled{\scriptsize b} anneau où $n$ inversible.

\[
 X_{\mathbb Z[1/n]}\simeq\widetilde R_{\mathrm{top}\,\mathbb Z[1/n]}
 \simeq\widehat R_{\mathbb Z\unsure{(1/n]|[1/n]}}
\]

% slanted
\margin{Sous-topos ouvert}

top $\mathbb Z[1/n]$, top. où $X_i\to X$ est couvrante ssi l'un des
$X_i\otimes_{\mathbb Z}\mathbb Z[1/n]\to\deleted{X}X\otimes_{\mathbb Z}\mathbb Z[1/n]$ a une section

$R_{\mathbb Z[1/n]}$, sous-cat. pleine de $R$ des $\mathbb Z[1/n]$ schémas \inserted{aff.} de t.f.

\deleted{\textcircled{\scriptsize c} Soit plus généralement $S\to e_{\widehat R}$ un mono dans $\mathrm{Pro}\,R$, $R=\text{Sch-aff}$ \inserted{$\mathbb Z\to\Lambda$ (épi d'ann)}
d'où $R_{/S}\hookrightarrow R$ sous-catégorie pleine de $R$ (corresp. aux \deleted{\unsure{annx}} anneaux de t.f. sur $\mathbb Z$ qui sont des $\Lambda$-algèbres).}

\deleted{Considérons les Anneaux du Topos $\mathcal E$ qui sont $\Lambda$-algèbres.
(La structure de $\Lambda$ algèbre est elle unique ? Oui, car $\Lambda$ est commençons par le procédé \unsure{standart|standard} au cas ensembliste).
les $\Lambda$-algèbres de $\mathcal E$ sont les fonct. \unsure{foncteurs|ponctuels} $\mathcal E\to\deleted{\rm Ens}(\Lambda\text{-alg.}\leftrightarrow\mathrm{Ann})$
qui sont représentables ensemblistement. Or \deleted{\unsure{il existe}}
\deleted{\unsure{($\Lambda$-alg) $\simeq\operatorname{Hom}$}} On identifiera (Ann) à $\operatorname{Hom}$ (\unsure{$R$, Ens})\unsure{$^{\rm opp}$}}
